\documentclass[11pt]{article}

\usepackage[a4paper,margin=28mm]{geometry}
\usepackage{amsmath,amssymb,amsthm,mathtools}
\usepackage{microtype}
\usepackage[hidelinks]{hyperref}

\newtheorem{theorem}{Theorem}
\newtheorem{proposition}[theorem]{Proposition}
\newtheorem{corollary}[theorem]{Corollary}
\theoremstyle{remark}
\newtheorem{remark}[theorem]{Remark}

\newcommand{\Z}{\mathbb Z}
\newcommand{\Q}{\mathbb Q}
\newcommand{\A}{\mathbb A}

\title{A Polynomial Parametrization of\texorpdfstring{\\[2mm]
\(xyz+t^2-1=0\)}{ xyz+t^2-1=0}}
\author{}
\date{}

\begin{document}

\maketitle
\vspace{-2.2em}

\begin{abstract}
We exhibit two polynomial maps over \(\Z\) whose union parametrizes every
rational solution of \(xyz+t^2-1=0\).  Integer parameter values yield an
explicit three-parameter family of integer solutions.  We also identify the
inverse map on the nondegenerate locus, record its exact integrality criterion,
and explain the factorization that leads naturally to the construction.
\end{abstract}

\section{The parametrization}

Let
\[
X=\bigl\{(x,y,z,t): xyz+t^2-1=0\bigr\}.
\]
For each \(\varepsilon\in\{1,-1\}\), define
\[
\Phi_{\varepsilon}(a,b,c)
 =\bigl(a,\ b,\ c(2-abc),\ \varepsilon(1-abc)\bigr).
\]
Every coordinate of \(\Phi_\varepsilon\) is a polynomial in \(a,b,c\)
with integer coefficients.

\begin{theorem}\label{thm:main}
The following statements hold.
\begin{enumerate}
    \item For every commutative ring \(R\) and every \((a,b,c)\in R^3\),
    the point \(\Phi_\varepsilon(a,b,c)\) satisfies
    \(xyz+t^2-1=0\).

    \item The two polynomial maps are jointly exhaustive over the rationals:
    \[
       X(\Q)=\Phi_{1}(\Q^3)\cup\Phi_{-1}(\Q^3).
    \]

    \item On the open locus \(xy\ne0\), the map \(\Phi_1\) has the inverse
    \[
       (x,y,z,t)\longmapsto
       \left(x,\ y,\ \frac{1-t}{xy}\right).
    \]
    Consequently, \(X\) is a rational affine threefold; more precisely,
    \(\Phi_1\) identifies the open set \(ab\ne0\) in \(\A^3_\Q\) with the
    open set \(xy\ne0\) in \(X\).
\end{enumerate}
\end{theorem}

\begin{proof}
Put \(s=abc\).  For
\((x,y,z,t)=\Phi_\varepsilon(a,b,c)\), one has
\[
xyz=ab\,c(2-abc)=s(2-s),
\]
and, since \(\varepsilon^2=1\),
\[
t^2-1=(1-s)^2-1=s^2-2s=-s(2-s).
\]
Their sum is zero.  This proves the first assertion as a polynomial
identity over \(\Z\), and hence over every commutative ring.

Now let \((x,y,z,t)\in X(\Q)\).  Suppose first that \(xy\ne0\), and set
\[
a=x,\qquad b=y,\qquad c=\frac{1-t}{xy}.
\]
Then \(abc=1-t\), so the fourth coordinate of \(\Phi_1(a,b,c)\) is \(t\).
Moreover,
\[
\begin{aligned}
c(2-abc)
 &=\frac{1-t}{xy}\bigl(2-(1-t)\bigr)\\
 &=\frac{(1-t)(1+t)}{xy}\\
 &=\frac{1-t^2}{xy}=z,
\end{aligned}
\]
where the final equality uses \(xyz=1-t^2\).  Thus
\((x,y,z,t)=\Phi_1(a,b,c)\).

If \(xy=0\), then the defining equation gives \(t^2=1\).  When \(t=1\),
\[
(x,y,z,t)=\Phi_1\left(x,y,\frac z2\right),
\]
and when \(t=-1\),
\[
(x,y,z,t)=\Phi_{-1}\left(x,y,\frac z2\right).
\]
This proves the second assertion.  On \(xy\ne0\), the preceding formula
for \(c\) is forced by \(t=1-abc\), so it is the inverse of \(\Phi_1\).
The final statement follows.
\end{proof}

\begin{corollary}\label{cor:integers}
For arbitrary \(a,b,c\in\Z\),
\[
\boxed{
(x,y,z,t)=\bigl(a,b,c(2-abc),\,\pm(1-abc)\bigr)
}
\]
is an integer solution.  In particular, the equation has infinitely many
integer solutions.
\end{corollary}

\begin{proof}
The first statement is Theorem~\ref{thm:main}(1) with \(R=\Z\).  For an
explicit one-parameter subfamily, take \(a=b=1\) and \(c=n\).  This gives
\[
(x,y,z,t)=\bigl(1,1,n(2-n),1-n\bigr),\qquad n\in\Z.
\]
Different values of \(n\) give different fourth coordinates, and hence
different solutions.
\end{proof}

\section{How the construction is found}

The equation is most naturally rewritten as
\[
xyz=1-t^2=(1-t)(1+t).
\]
The right-hand side already consists of two factors whose difference is
exactly \(2\).  To introduce three free parameters while preserving this
factorization, prescribe one of those factors to be their product:
\[
1-t=abc.
\]
This immediately forces
\[
t=1-abc,
\qquad
1+t=2-abc.
\]
Therefore
\[
1-t^2=(1-t)(1+t)=abc(2-abc).
\]
The product on the right can now be distributed among \(x,y,z\) by taking
\[
x=a,\qquad y=b,\qquad z=c(2-abc).
\]
This gives the displayed polynomial parametrization without any division.
Choosing instead the sign-reversed fourth coordinate produces the second
chart \(\Phi_{-1}\).

The essential idea is therefore simple but structural: force one linear
factor of \(1-t^2\) to absorb the product of the free parameters; the other
linear factor then remains polynomial because the two factors differ by
\(2\).

\section{Rational completeness versus integral completeness}

Theorem~\ref{thm:main} is exhaustive when the parameters are allowed to be
rational.  With integer parameters, the image is still a large
three-parameter family, but it does not contain every integer point.  The
precise obstruction on the nondegenerate locus is divisibility.

\begin{proposition}\label{prop:integrality}
Let \((x,y,z,t)\in X(\Z)\) with \(xy\ne0\).
\begin{enumerate}
    \item It has an integral preimage under \(\Phi_1\) if and only if
    \(xy\mid 1-t\).  In that case the preimage is uniquely
    \[
       \left(x,y,\frac{1-t}{xy}\right).
    \]

    \item It has an integral preimage under \(\Phi_{-1}\) if and only if
    \(xy\mid 1+t\).  In that case the preimage is uniquely
    \[
       \left(x,y,\frac{1+t}{xy}\right).
    \]
\end{enumerate}
\end{proposition}

\begin{proof}
For \(\Phi_1\), the first two coordinates force \(a=x\) and \(b=y\), while
the fourth coordinate forces
\(c=(1-t)/(xy)\).  Thus an integral preimage exists exactly when this
quotient is integral, and it is then unique.  The argument for
\(\Phi_{-1}\), whose fourth coordinate is \(abc-1\), is identical.
\end{proof}

\begin{remark}\label{rem:scope}
The distinction is genuine.  For example,
\[
(30,42,-4,71)\in X(\Z),
\]
because
\[
30\cdot42\cdot(-4)+71^2-1=-5040+5041-1=0.
\]
The pairwise products of the first three coordinates are
\(1260,-120,-168\), and none divides either
\(1-71=-70\) or \(1+71=72\).  Hence this point is not obtained from
integer parameters in either chart, even after permuting \(x,y,z\).
Thus the formulas give a complete polynomial parametrization over \(\Q\)
and an explicit infinite polynomial family over \(\Z\), but not a claimed
complete parametrization of all integral points by unrestricted integer
parameters.
\end{remark}

\end{document}
